\part{外观、标签与评价规则}


\section{主题与内置颜色}
\key{theme} 一次设置填充方式、渐变两端、轨道和文字颜色。
主题参数之后仍可逐项覆盖颜色；同一选项列表从左向右执行。
示例 29 集中展示 19 种内置配色，包括十种主题和九种单色。
渐变预设通过 \key{theme} 调用，内置颜色通过 \key{color} 设置纯色填充。
九种预定义单色为 \manualColor{lightgreen}、\manualColor{lightyellow}、
\manualColor{lightblue}、\manualColor{lightred}、\manualColor{rose}、
\manualColor{skyblue}、\manualColor{gold}、\manualColor{lavender} 和 \manualColor{peach}。

\begin{example}[before upper={\fontsize{9}{10.6}\selectfont\setlength{\parskip}{13.8pt}}]{示例 29：内置主题与颜色}
\gradbar[theme=blue]{72}

\gradbar[theme=teal]{72}

\gradbar[theme=solid]{72}

\gradbar[theme=gray]{72}

\gradbar[theme=lbyellow]{72}

\gradbar[theme=viblue]{72}

\gradbar[theme=cyblu]{72}

\gradbar[theme=violet]{72}

\gradbar[theme=orange]{72}

\gradbar[theme=forest]{72}

\gradbar[color=lightgreen]{72}

\gradbar[color=lightyellow]{72}

\gradbar[color=lightblue]{72}

\gradbar[color=lightred]{72}

\gradbar[color=rose]{72}

\gradbar[color=skyblue]{72}

\gradbar[color=gold]{72}

\gradbar[color=lavender]{72}

\gradbar[color=peach]{72}
\end{example}

\clearpage
\begingroup
\small
\renewcommand{\arraystretch}{1.15}
\begin{center}
\begin{tabularx}{\linewidth}{@{}c>{\centering\arraybackslash}X@{}}
\toprule
\multicolumn{1}{c}{\textbf{主题或颜色}} & \multicolumn{1}{c}{\textbf{特征与适用场景}}\\
\midrule
\key{blue} & 默认蓝色渐变，适合通用表格。\\
\key{teal} & 青绿色渐变，适合区分另一组指标。\\
\key{solid} & 蓝色纯色填充，减少渐变带来的视觉干扰。\\
\key{gray} & 灰度填充，适合黑白输出；仍应保留数值标签。\\
\midrule
\key{lbyellow} & 浅蓝至浅黄渐变，适合浅色背景下的柔和展示。\\
\key{viblue} & 浅紫至浅蓝渐变，适合低饱和度的对比表格。\\
\key{cyblu} & 青色至浅蓝渐变，适合突出一组重点指标。\\
\key{violet} & 紫罗兰至浅紫渐变，搭配淡紫轨道，适合区分另一组实验或指标。\\
\key{orange} & 暖橙至浅橙渐变，搭配淡橙轨道，适合突出重点结果。\\
\key{forest} & 森林绿至浅绿渐变，搭配淡绿轨道，适合完成度与资源指标。\\
\midrule
\key{lightgreen} & 浅绿纯色，适合展示完成度或达标指标。\\
\key{lightyellow} & 浅黄纯色，适合轻量提示；浅色背景下应保留深色标签。\\
\key{lightblue} & 浅蓝纯色，适合通用数据和大面积表格。\\
\key{lightred} & 浅红纯色，适合标出需要关注的指标。\\
\key{rose} & 玫红纯色，适合强调重点数据或区分对照组。\\
\key{skyblue} & 天蓝纯色，适合清晰呈现进度和比例。\\
\key{gold} & 金黄纯色，适合突出关键结果或优选方案。\\
\key{lavender} & 淡紫纯色，适合柔和区分不同类别。\\
\key{peach} & 蜜桃纯色，适合暖色风格的辅助指标。\\
\bottomrule
\end{tabularx}
\end{center}
\endgroup

\begin{example}{示例 30：自定义渐变与轨道}
\gradbar[
  theme=blue,
  left color=violet!80!blue,
  right color=violet!20,
  track color=violet!5
]{68}

\gradbar[
  theme=teal,
  fill=solid,
  left color=orange!85!black,
  track color=orange!10
]{68}
\end{example}

\subsection{颜色参数}
内置配色无需自行定义。主题和渐变预设使用 \key{theme}；单色使用 \key{color}。
单色名称也可用于 \key{left color} 或 \key{right color}，名称不带反斜杠。

\key{theme=blue,color=rose} 使用玫红纯色；
\key{color=rose,theme=blue} 则恢复蓝色渐变。

\key{left color} 是渐变起点，也决定纯色填充的颜色；
\key{right color} 仅在 \key{fill=gradient} 时生效。
\key{track color} 控制空轨道，\key{text color} 控制标签。
颜色可使用 xcolor 的混色表达式。

\note{不要只靠颜色传递关键含义。并列指标应在列标题中写明名称和单位，
并保留可读的数值。可设置 \key{text color=auto} 根据背景选择黑白文字，也可显式指定颜色。}

\subsection{色觉友好的颜色角色与冗余线索}
\key{palette=accessible} 提供七种 Okabe--Ito 颜色角色与循环堆叠纹理；
\key{theme=accessible} 选用对应实心条形角色。
Heat 可用 \key{heat palette} (\key{accessible-sequential}, \key{accessible-diverging})，
原生表格可用 \key{theme=accessible} 和 \key{style=accessible}，旧默认值保持不变。
\begin{example}[text and listing,sidebyside=false]{颜色、文字、纹理与打印}
\gradstack[palette=accessible,width=45mm,height=16pt,
  stack names={A,B,C},segment labels=name]{30,20,50}
\par\medskip
\gradstack[palette=accessible,print,width=45mm,height=16pt,
  stack names={A,B,C},segment labels=name]{30,20,50}
\end{example}
颜色与纹理会循环，超过七类时尤其需要保留名称与不同符号。
渐变标签的自动对比度只是中心估计；必要时采用实心填充、外置标签或白底标签。
Heat 按实际底色选择黑白文字。缺失标记、空心异常点和打印纹理提供颜色之外的线索。
替代文本保留原值、单位和统计口径；长描述可由显式 \key{alt} 替换，并在邻近文字提供完整数据。
配色模拟与结构检查不构成普遍可访问性认证。
